\documentclass[a4paper, fontsize=11pt]{article}
\usepackage[T2A]{fontenc}
\usepackage[unicode, pdftex]{hyperref}
\usepackage{amsmath}
\usepackage{amssymb}

\usepackage{geometry}
\geometry{margin=1cm}

\begin{document}



\normalsize

\noindent
{\LARGE \textbf{Igor Ignashin}}\\[0.5em]

\noindent
\textbf{Email:} \href{mailto:ignashin.in@phystech.edu}{ignashin.in@phystech.edu} \\
\textbf{GitHub:} \href{https://github.com/ThunderstormXX}{github.com/ThunderstormXX}

\vspace{0.5cm}

\textbf{Profile:}
\begin{itemize}
    \item Research engineer focused on LLM post-training, SFT, inference efficiency, and optimization. I build reproducible training and evaluation pipelines for teacher-generated SFT data, reward parsing, RLVR/GRPO-style math training, benchmark reporting, and vLLM-based early-exit/adaptive decoding.
    \item Strong mathematical background in optimization and stochastic processes, with research experience connecting theory, experiments, and practical large-model training workflows.
\end{itemize}

\textbf{Technical Skills:}
\begin{itemize}
    \item Languages: Python, C++, SQL
    \item LLM Training/Post-Training: SFT, teacher distillation, RLVR/GRPO-style training, reward functions and graders, synthetic and teacher-generated data, evaluation pipelines
    \item ML/Optimization: PyTorch, JAX, vLLM, TRL, Hugging Face workflows, RL, LLMs, convex and nonconvex optimization, game theory
    \item Mathematics: stochastic processes, optimization, minimax problems, probability theory
    \item Data/Infra: multi-GPU server workflows, JSONL data pipelines, YQL, DataLens, Nirvana, Linux, Git
\end{itemize}


\textbf{Work Experience:}
\begin{itemize}
    \item Visiting Research Student, MBZUAI (under \href{https://scholar.google.com/citations?user=QPVriwoAAAAJ&hl=ru&oi=ao}{Eduard Gorbunov}) \hfill Jan 2026 -- Mar 2026
\begin{itemize}
    \item Conducted theoretical research on LMO-based optimization methods and acceleration schemes, including analyses of the Muon optimizer and Lookahead-type techniques.
\end{itemize}

    \item Laboratory of Mathematical Optimization Methods, MIPT / ISP RAS / BRAIn Lab \hfill 2024--Present
    \begin{itemize}
        \item Developing LLM post-training and inference-efficiency projects, including a Qwen/DeepScaleR math pipeline with teacher-generated SFT data, reward parsing, RLVR/GRPO-style training, and benchmark reporting, and a vLLM-based early-exit/adaptive-decoding evaluation pipeline.

        \item Working on a Huawei project on neural network-based schedulers, with a focus on reducing the performance gap between classical and learned approaches.
        \item Proved local convergence of the ALSO method for nonconvex and $\mu$-strongly concave minimax problems in \textit{``Aligning Distributionally Robust Optimization with Practical Deep Learning Needs''}.

        \item Leading student research projects at YSDA and MIPT AI360 on LLM inference acceleration, distillation, multi-agent RL, game-theoretic LLM training, SGD dynamics, and applied optimization, with an emphasis on theory, experiments, and publication-quality research.

        % \item Leading a student project at Yandex School of Data Analysis (YSDA) on accelerating LLM inference via early exiting (2025).

        % \item Leading a project on neural network agents in game-theoretic settings for 2nd–3rd year students of MIPT (AI360 program)(2025).

        % \item Conducting research on LLM distillation, pruning, and interpretability (2025).

        \item Collaborating with \href{https://scholar.google.com/citations?user=RNnvD3wAAAAJ\&hl=en}{Andrei Leonidov's group} on SGD analysis and multi-agent reinforcement learning, including SGD noise analysis, theoretical frameworks for nonconvex optimization dynamics, and joint research on multi-agent LLM systems in game-theoretic scenarios and social dilemmas.

        % \item Collaborating with colleagues from Tinkoff Bank on training strategies for LLM agents. Currently exploring modifications of the DPO loss, which implicitly incorporates transitive rewards (introducing a partial order on model responses). Investigating multi-agent reward mechanisms such as alpha-rank for handling non-transitive comparisons.
        % \item Worked on optimization of DL models for tabular data, including experiments with the \href{https://github.com/yandex-research/tabm}{framework by Yandex Research}. Supervised a student project on improving periodic embeddings for tabular tasks (2024).
        \item Researched equilibrium traffic assignment problems and Frank-Wolfe algorithm modifications, resulting in two journal publications and two conference presentations. Currently developing stochastic extensions of the LMO oracle and exploring applications to game-theoretic formulations.

    \end{itemize}

    \item Yandex, Data Analyst Intern \hfill June 2022 -- September 2022
    \begin{itemize}
        \item Automated a daily pipeline for ticket distribution and closure analytics for B2B contractors, improving transparency and monitoring efficiency.
        \item Built data processing workflows in YQL and designed dashboards in DataLens for operational insights.
        \item Worked with Yandex internal tools, including Nirvana, for orchestrating data workflows.
    \end{itemize}
\end{itemize}

\textbf{Selected LLM Projects:}
\begin{itemize}
    \item \textbf{Qwen math SFT/RLVR pipeline:} set up a post-training workflow from base evaluation to SFT, RLVR/GRPO-style training, final evaluation, and Base vs SFT vs SFT+RL reporting for Qwen models.
    \item \textbf{DeepScaleR teacher-SFT distillation:} built a pipeline for generating teacher solutions, scoring them with local math reward/parsing code, training a student model on SFT JSONL, and evaluating on AIME 2024, MATH 500, AMC 2023, Minerva Math, and OlympiadBench.
    \item \textbf{vLLM early-exit/adaptive decoding:} built a reproducible offline inference and evaluation contour with deterministic sharding, resumable JSONL outputs, merge/progress scripts, metrics, and AdaDecode baseline integration on multi-GPU servers.
    \item \textbf{LLM pruning and healing:} worked on layer pruning, LoRA healing, decoder search, and attention-head merging experiments for LLM compression.
\end{itemize}


\textbf{Selected Publications:}

\textbf{\href{https://scholar.google.com/citations?hl=ru&pli=1&user=YXllYbsAAAAJ}{My Google Scholar}}


\begin{itemize}
    \item \href{https://arxiv.org/pdf/2508.16734}{Aligning Distributionally Robust Optimization with Practical Deep Learning Needs}. Accepted to a NeurIPS 2025 workshop.

    \item NeurIPS 2026 submission on stochastic dynamics of SGD, proposing a new perspective beyond Brownian-motion assumptions
    (\href{https://openreview.net/forum?id=qqzq4YOnAo&referrer=%5Bthe%20profile%20of%20Igor%20Ignashin%5D(%2Fprofile%3Fid%3D~Igor_Ignashin1)}{\textit{Why SGD is not Brownian Motion: A New Perspective on Stochastic Dynamics}}).

    % \item Submitted two papers to ICLR 2026: one on SGD noise modeling (\href{https://openreview.net/forum?id=ToATXZ0E0f&referrer=%5Bthe+profile+of+Igor+Ignashin%5D%28%2Fprofile%3Fid%3D%7EIgor_Ignashin1%29}{\textit{Why SGD is not Brownian Motion: A New Perspective on Stochastic Dynamics}}), and another on convergence of ALSO in nonconvex–strongly concave settings (\href{https://openreview.net/forum?id=p28VywxRP1&referrer=%5Bthe%20profile%20of%20Igor%20Ignashin%5D(%2Fprofile%3Fid%3D~Igor_Ignashin1)}{\textit{Aligning Distributionally Robust Optimization with Practical Deep Learning Needs}}).

    % \item Submitted a paper on LLM distillation and interpretability to a workshop. Proposed an iterative layer-dropping method for decoder distillation with LoRA fine-tuning. \href{https://openreview.net/forum?id=KfuiwhQcJH\&referrer=\%5Bthe\%20profile\%20of\%20Igor\%20Ignashin\%5D(\%2Fprofile\%3Fid\%3D~Igor\_Ignashin1)}{The Ridiculous Redundancy of LLMs: Which Layers and Heads Actually Matter}

    \item \href{https://openreview.net/forum?id=KfuiwhQcJH}{The Ridiculous Redundancy of LLMs: Which Layers and Heads Actually Matter}. Work on LLM layer pruning, LoRA healing, and attention-head merging for decoder compression.

    \item \href{http://crm.ics.org.ru/journal/article/3433/}{Publication} in \textit{Computer Research and Modeling}: \href{https://crm.ics.org.ru/uploads/crmissues/crm_2024_1/14_ignashin.pdf}{Frank-Wolfe Modifications for Equilibrium Traffic Assignment}.

    \item Preprint: \href{https://arxiv.org/pdf/2509.23840}{Stochastic Origin Frank-Wolfe for Traffic Assignment}. Presented at \href{https://siriusmathcenter.ru/056w}{TFN-2025}. Accepted to the \href{https://mathbooks.ru/Springer-Sirius/}{Journal of Mathematical Sciences, Series B}.

    \item Preprint: \href{https://arxiv.org/abs/2509.13392}{Modeling Skier Flows via Wardrop Equilibrium in Closed Capacitated Networks}. Presented at \href{https://siriusmathcenter.ru/056w}{TFN-2025}. Accepted to the \href{https://mathbooks.ru/Springer-Sirius/}{Journal of Mathematical Sciences, Series B}.

    % \item \href{https://old.mipt.ru/upload/medialibrary/d79/fpmi.pdf}{Thesis} “On equilibrium traffic assignment problem” at the 66th All-Russian Scientific Conference of MIPT.

    \item \textit{Conjugate Frank-Wolfe in Machine Learning.} Presented at OPTIMA and accepted to CCIS.

\end{itemize}

\textbf{Education:}
\begin{itemize}
    \item B.S. in Applied Mathematics and Physics, Moscow Institute of Physics and Technology (MIPT), Faculty of Applied Mathematics and Informatics.
    \item M.S. candidate, MIPT, Faculty of Applied Mathematics and Informatics, Department of Intelligent Data Analysis.
\end{itemize}

\textbf{Summer Schools and Hackathons:}
\begin{itemize}
    \item AIRI Institute Summer School on Artificial Intelligence 2025
    \item Summer School “Control, Information and Optimization” in memory of B. T. Polyak 2025
    \item Summer School and Hackathon in Financial Mathematics, Quantathon 2024
    \item 5th and 6th International Summer Schools on Financial Mathematics (2024--2025)
\end{itemize}

\textbf{School Olympiad Achievements:}
\begin{itemize}
    \item Winner, Phystech Mathematics Olympiad.
    \item Winner, Step into the Future Olympiads in Mathematics and Physics.
    \item Winner, KFU Interregional Mathematics Olympiad.
    \item Winner, municipal stage of the All-Russian School Olympiad in Mathematics.
\end{itemize}



\end{document}
